% !TEX root = ../main.tex

\chapter{Foundations of PDEs}

\section{Introduction and Notation}

A partial differential equation (PDE) is an equation involving an unknown function of several variables and its partial derivatives. A typical \(k\)-th order equation can be written in the compact form
\[
  F\bigl(x,u,Du,D^2u,\ldots,D^k u\bigr)=0,
\]
where \(x\in \R^n\), \(u=u(x)\), and \(D^j u\) denotes all derivatives of order \(j\).

\begin{definitionbox}{Second-Order Linear PDE}
Let \(u=u(x)\), \(x\in \Omega\subset\R^n\). A second-order linear PDE has the form
\[
  \sum_{i,j=1}^{n} a_{ij}(x)\frac{\partial^2 u}{\partial x_i\partial x_j}
  + \sum_{i=1}^{n} b_i(x)\frac{\partial u}{\partial x_i}
  + c(x)u = f(x).
\]
When \(f=0\), the equation is homogeneous; otherwise it is nonhomogeneous.
\end{definitionbox}

The coefficient matrix \(A(x)=(a_{ij}(x))\) encodes the highest-order part of the equation. Much of the qualitative behavior of solutions is already visible in the signature of this matrix.

\begin{definitionbox}{Classification by Principal Part}
At a point \(x_0\in\Omega\), a second-order equation is commonly classified as follows:
\begin{itemize}[leftmargin=2em]
  \item elliptic if the eigenvalues of \(A(x_0)\) have the same sign;
  \item hyperbolic if exactly one eigenvalue has the opposite sign;
  \item parabolic if \(A(x_0)\) is degenerate, for example \(\det A(x_0)=0\).
\end{itemize}
\end{definitionbox}

\begin{examplebox}{Classical Models}
Three standard equations on a domain in \(\R^n\) are:
\begin{enumerate}[label=(\roman*),leftmargin=2.4em]
  \item the Poisson equation, \(-\Delta u=f\);
  \item the wave equation, \(u_{tt}-c^2\Delta u=0\);
  \item the heat equation, \(u_t-\kappa\Delta u=0\).
\end{enumerate}
\end{examplebox}

\section{Initial and Boundary Data}

PDEs usually have infinitely many formal solutions until extra data are prescribed. For evolution equations one often gives initial data at \(t=0\). For boundary value problems one prescribes information along \(\partial\Omega\).

\begin{definitionbox}{Cauchy Problem}
Given a hypersurface \(\Gamma\), a Cauchy problem asks for a solution near \(\Gamma\) whose value and selected normal derivatives are prescribed on \(\Gamma\).
\end{definitionbox}

For equations on a bounded spatial region, the most common boundary conditions are
\begin{itemize}[leftmargin=2.0em]
  \item \textbf{Dirichlet:} \(u|_{\partial\Omega}=g\);
  \item \textbf{Neumann:} \(\frac{\partial u}{\partial n}\big|_{\partial\Omega}=h\);
  \item \textbf{Robin:} \(\left(\frac{\partial u}{\partial n}+\alpha u\right)\big|_{\partial\Omega}=k\).
\end{itemize}

\begin{principlebox}{Well-Posedness (Hadamard)}
A PDE problem is called well posed if it satisfies:
\begin{enumerate}[label=(\alph*),leftmargin=2.2em]
  \item existence: at least one solution exists;
  \item uniqueness: at most one solution exists;
  \item stability: the solution depends continuously on the prescribed data.
\end{enumerate}
\end{principlebox}

\section{Well-Posedness Framework}

Hadamard's criteria form a basic test for a mathematical model. A problem that is not stable may fit data exactly and still be useless for prediction, because small measurement errors can create large changes in the solution.

\begin{notebox}{Note}
Modern PDE theory often proves stability in weak norms. For example, the variational treatment of elliptic equations uses Hilbert spaces and coercive bilinear forms rather than pointwise second derivatives.
\end{notebox}

\section{Superposition and Duhamel Principle}

For a linear operator \(L\), the equation \(Lu=f\) separates the geometry of the operator from the forcing term. If \(u_1\) and \(u_2\) solve homogeneous equations, then any linear combination is again a homogeneous solution. For time-dependent equations, Duhamel's principle rewrites a forced solution as the accumulation of homogeneous responses.
